\documentclass[letterpaper,11pt]{article}

\input{shared/preamble}
\input{shared/commands}

\begin{document}

\begin{center}
  \resumeName{Harinandan Kotamsetti}
    \small Cupertino, CA $|$ +1 (669) 264-9458 $|$
    \href{mailto:hareee234@gmail.com}{\underline{hareee234@gmail.com}} $|$
    \href{https://www.linkedin.com/in/hareee234/}{\underline{linkedin.com/in/hareee234}} $|$
    \href{https://github.com/AwesomenessReborn}{\underline{github.com/AwesomenessReborn}}
\end{center}

\section{Education}
  \resumeSubHeadingListStart
    \resumeSubheading
      {\textbf{San Jose State University}}{San Jose, CA}
      {Bachelor of Science in Computer Engineering}{Expected Dec 2026}
      \resumeItemListStart
        \resumeItem{\textbf{Coursework:} OS, Computer Architecture, C/C++, FPGA, Verilog, Linear Algebra, x86 \& MIPS Assembly}
      \resumeItemListEnd
  \resumeSubHeadingListEnd

\section{Experience}
  \resumeSubHeadingListStart
    \resumeSubheading
      {\textbf{STMicroelectronics}}{Santa Clara, CA}
      {\textbf{Software \& Algorithm Intern}}{Mar 2025 -- Present}
      \resumeItemListStart
        
        \resumeItem{
          Designed a \textbf{103K-parameter CNN-LSTM} for walking-speed estimation, reaching \textbf{0.059 m/s MAE} on walking logs; quantized it to \textbf{INT8} \& deployed on \textbf{Snapdragon NPU via Qualcomm QNN}. Demoed at \textbf{CES 2026}.
        }

        \resumeItem{
          Implemented a \textbf{geodesic angle loss and adaptive loss-weighting changes} for VR head-motion forecasting, cutting error \textbf{9\% at 111 ms} and \textbf{41\% at 11 ms}; shortened input context from \textbf{3.3 s to 1.1 s} with \textbf{no measurable accuracy loss}, speeding training \textbf{$\sim$3$\times$}.
        }
        
        \resumeItem{
          \textbf{Owned the head-motion model's training pipeline for runs consuming 45+ GPU-hours}, increasing throughput \textbf{2.5$\times$} with bf16 + \texttt{torch.compile} optimizations and implementing per-horizon best-checkpoint saving after discovering $\sim$79\% of evaluated checkpoints were not retained.
        }
        
        \resumeItem{
          Ported an \textbf{887K-parameter VR head-motion model} to \textbf{Qualcomm's Hexagon NPU} (PyTorch $\rightarrow$ ONNX $\rightarrow$ QAIRT), matching PyTorch output within \textbf{0.056$^\circ$ on-device}.
        }
        
        \resumeItem{
          \textbf{Deployed always-on fall detection and wrist-tilt recognition entirely in-sensor} (ISPU, zero host-MCU compute); profiled 104 Hz on-sensor fusion at \textbf{3.08 ms worst case} (32\% of budget).
        }

      \resumeItemListEnd

    \resumeSubheading
      {\textbf{Spartan Racing (Formula SAE)}}{San Jose, CA}
      {\textbf{Embedded Systems \& Software Developer}}{Aug 2023 -- Feb 2025}
      \resumeItemListStart
        
        \resumeItem{
          \textbf{Improved real-time BMS observability} by building a \textbf{CAN --- InfluxDB --- Grafana} telemetry pipeline for live diagnostics, fault visibility, and team-wide debugging across a \textbf{400 V Formula SAE} battery system.
        }

        \resumeItem{
          \textbf{Cut STM32 BMS polling latency 14$\times$ (7.0 s $\rightarrow$ 0.5 s)} by eliminating blocking CAN/thermistor delays and adding mailbox-aware CAN transmit logic.
        }
    
        \resumeItem{
          \textbf{Implemented state-of-charge estimation} (Coulomb counting with open-circuit-voltage lookup-table correction) for a \textbf{400 V, 768-cell BMS on STM32}, reaching \textbf{3\% MAE}.
        }
      
      \resumeItemListEnd
  \resumeSubHeadingListEnd

%-----------PATENTS & INVENTION DISCLOSURES-----------
\section{Patents \& Invention Disclosures}
  \resumeSubHeadingListStart
    \resumePatentHeading
      {Glasses with Wellness and Health Monitoring}
      {STMicroelectronics, Co-inventor}
      {Filed May 2026}
    \resumePatentHeading
      {Motion-to-Photon Latency Compensation for XR/AR/VR Glasses}
      {STMicroelectronics, Co-inventor}
      {Invention Disclosure, 2026}
  \resumeSubHeadingListEnd

\section{Projects}
  \resumeSubHeadingListStart

    \resumeProjectHeading
        {\textbf{Nemotron-H Fine-Tuning --- Kaggle NLP Competition} $|$
          \emph{Python, PyTorch, LoRA, llama.cpp}}{2026}
        \resumeItemListStart
            \resumeItem{
              \textbf{Improved controlled task accuracy by +19.6-22.5 absolute percentage points} by building a LoRA fine-tuning and evaluation pipeline for Nemotron-H 4B (Hybrid Mamba-Attention), \textbf{training 20.2M / 3.99B parameters (0.5\%) on 8,000+ self-generated reasoning traces} with assistant-only label masking.
            }
            \resumeItem{
              \textbf{Made Nemotron-H fine-tuning stable} after the stock Hugging Face Trainer and \texttt{generate()} diverged: traced an FP32/BF16 mismatch in the Mamba kernels, restricted LoRA to attention/MLP layers, and wrote a custom loop with no-cache chunked inference to avoid corrupted SSM state.
            }
        \resumeItemListEnd

    \resumeProjectHeading
        {\textbf{MERIT UAV --- Autonomous Surveillance Drone} $|$
          \emph{Python, YOLOv8n, CUDA, Jetson Orin, MAVLink}}{2026}
        \resumeItemListStart
            \resumeItem{
              \textbf{Ran YOLOv8n detection on-board a Jetson Orin Nano at 30 FPS}, geotagging each detection with live GPS and uploading it during flight-tested autonomous missions.
            }
            \resumeItem{
              \textbf{Implemented and flight-tested autonomous MAVLink control} (takeoff, altitude, position, yaw) from the Jetson to a Pixhawk, gating detections on GPS lock so every upload carries a valid location.
            }
        \resumeItemListEnd

  \resumeSubHeadingListEnd

\section{Technical Skills}
  \resumeSubHeadingListStart
    \small{
      \item{
        \textbf{Languages}{: C/C++, Embedded C Programming,  Python, Kotlin, Java, SQL} \\
        \textbf{ML \& AI}{: PyTorch, Tensorflow, TFLite, ONNX, QNN/QAIRT, LoRA fine-tuning, scikit-learn, NumPy, Pandas, SciPy.} \\
        \textbf{Signal Processing \& Analysis}{: MATLAB (familiar), Simulations, analytical modeling, feasibility studies, Digital filtering (Butterworth), feature extraction (FFT, windowing), time-series analysis, sensor fusion.} \\
        \textbf{Developer Tools}{: Android (NDK/JNI, Gradle, adb/logcat), CMake, Docker, GitHub Actions, CI/CD, Claude Code, OpenCode (agentic coding workflows, code review augmentation).} \\
        \textbf{Data \& Infrastructure}{: PostgreSQL, InfluxDB, Grafana, AWS S3.} \\
        \textbf{Concepts}{: On-device ML inference, model deployment, real-time data pipelines, model evaluation, debugging.}
      }
    }
  \resumeSubHeadingListEnd

\end{document}